\documentclass[12pt,a4paper]{article}
\usepackage[margin=1in]{geometry}
\usepackage{graphicx}
\usepackage{hyperref}
\usepackage{xcolor}
\usepackage{booktabs}
\usepackage{tabularx}
\usepackage{enumitem}
\usepackage{fancyhdr}
\usepackage{titlesec}
\usepackage{tcolorbox}
\usepackage{listings}
\usepackage{fontawesome5}
\usepackage{amsmath}

\definecolor{primary}{HTML}{0A66C2}
\definecolor{accent}{HTML}{057642}
\definecolor{warning}{HTML}{B24020}
\definecolor{codebg}{HTML}{1E1E1E}
\definecolor{codetext}{HTML}{D4D4D4}
\definecolor{darktext}{HTML}{1D2226}

\hypersetup{colorlinks=true,linkcolor=primary,urlcolor=primary}
\pagestyle{fancy}
\fancyhf{}
\fancyhead[L]{\textcolor{primary}{\textbf{LinkedIn Comment Generator}}}
\fancyhead[R]{\textcolor{gray}{Technical Document v2.0}}
\fancyfoot[C]{\thepage}

\titleformat{\section}{\Large\bfseries\color{primary}}{}{0em}{}
\titleformat{\subsection}{\large\bfseries\color{darktext}}{}{0em}{}
\titleformat{\subsubsection}{\normalsize\bfseries\color{darktext}}{}{0em}{}

\lstset{
    backgroundcolor=\color{codebg},
    basicstyle=\ttfamily\small\color{codetext},
    breaklines=true,
    frame=single,
    rulecolor=\color{gray},
    tabsize=2,
    showstringspaces=false,
}

\newtcolorbox{archbox}[1]{colback=blue!3,colframe=primary,title=#1,fonttitle=\bfseries}

\begin{document}

% ============================================
% TITLE PAGE
% ============================================
\begin{titlepage}
\centering
\vspace*{2cm}
{\Huge\bfseries\color{primary} LinkedIn Comment Generator\par}
\vspace{0.5cm}
{\Large\color{darktext} Technical Architecture Document\par}
\vspace{1cm}
{\large Version 2.0 --- FastAPI + LangGraph\par}
\vspace{0.5cm}
{\large April 2026\par}
\vspace{2cm}

\begin{tcolorbox}[colback=blue!3,colframe=primary,width=0.85\textwidth]
\small
\begin{tabular}{ll}
\textbf{Frontend:} & Next.js 16 (App Router) + Tailwind CSS + TypeScript \\
\textbf{Backend:} & FastAPI + LangGraph (Python 3.12) \\
\textbf{AI Engine:} & Groq API (Llama 3.3 70B + Llama 4 Scout Vision) \\
\textbf{Search:} & SerpAPI (Google Search) \\
\textbf{Data:} & RapidAPI LinkedIn Scraper \\
\textbf{Database:} & SQLite \\
\textbf{RAG:} & 1,719 reference comments \\
\textbf{Hosting:} & Railway (single container) \\
\textbf{Repository:} & \texttt{github.com/lucasleadfreak-halfengi/linkedin-comment-generator} \\
\end{tabular}
\end{tcolorbox}

\vfill
\end{titlepage}

\tableofcontents
\newpage

% ============================================
% SYSTEM ARCHITECTURE
% ============================================
\section{System Architecture}

\subsection{High-Level Overview}

\begin{archbox}{Architecture Diagram}
\begin{verbatim}
Browser (Railway URL)
    │
    ▼
┌─────────────────────────────────────────────────┐
│           Single Railway Container               │
│                                                  │
│  Next.js Frontend (port 3000)                   │
│  ├── Upload Page        ──proxy──┐              │
│  ├── Dashboard (Feed)            │              │
│  ├── Post Detail                 │  /api/*      │
│  └── Settings                    ▼              │
│                          FastAPI Backend (8000)  │
│                          ├── LangGraph Pipeline  │
│                          ├── SQLite Database     │
│                          └── RAG (1719 examples) │
└─────────────────────────────────────────────────┘
         │              │              │
    ┌────┘        ┌─────┘        ┌─────┘
    ▼             ▼              ▼
 RapidAPI     Groq API       SerpAPI
 (Scrape)    (AI Engine)    (Search)
\end{verbatim}
\end{archbox}

\subsection{Why This Architecture}

\begin{center}
\begin{tabularx}{\textwidth}{lX}
\toprule
\textbf{Decision} & \textbf{Rationale} \\
\midrule
LangGraph over linear pipeline & Retry on low quality, parallel generation, conditional flow \\
FastAPI over Next.js API routes & Python ecosystem, reuse existing code, better AI tooling \\
Next.js proxy to FastAPI & Single deployment, no CORS issues, clean separation \\
SQLite over PostgreSQL & Zero config, single file, sufficient for this scale \\
Groq over OpenAI/Anthropic & 460 tok/s speed, \$0.59/M tokens (10x cheaper than GPT-4) \\
RAG over fine-tuning & Instant updates, no training cost, \$0 compute \\
\bottomrule
\end{tabularx}
\end{center}

% ============================================
% LANGGRAPH PIPELINE
% ============================================
\section{LangGraph Pipeline}

\subsection{Pipeline Graph}

\begin{archbox}{LangGraph Node Flow}
\begin{verbatim}
START
  │
  ▼
[analyze] ──── Groq llama-3.3-70b (JSON mode)
  │             Detects: post_type, post_tone, sentiment,
  │             specific_details, best_response_angles
  ▼
[search] ───── SerpAPI (conditional)
  │             Only for: advice, controversial, philosophical
  │             Skips: celebration, hiring, milestone
  ▼
[rag] ──────── SQLite query → top 5 examples
  │             Filters by post_type, sorts by engagement_score
  │             Extracts angle names from examples
  ▼
[generate] ─── 3x Groq calls IN PARALLEL
  │             Each with different angle + temperature
  │             Variation 1: temp=0.70
  │             Variation 2: temp=0.85
  │             Variation 3: temp=1.00
  ▼
[humanize_validate] ── String processing + quality scoring
  │
  ▼
quality < 50? ──YES──▶ [generate] (retry, max 2x)
  │
  NO
  ▼
END ── 3 comments saved to DB
\end{verbatim}
\end{archbox}

\subsection{State Schema}

\begin{lstlisting}
class PipelineState(TypedDict):
    post_id: int
    post_text: str
    post_author: dict           # {name, headline}
    media: list[dict]
    analysis: dict              # PostAnalysis
    web_context: str            # Formatted search results
    rag_context: str            # Formatted RAG examples
    rag_angles: list[str]       # Angle names from RAG
    profile: dict               # CommenterProfile
    comments: Annotated[list[dict], operator.add]  # APPENDS
    retry_count: int
    errors: Annotated[list[str], operator.add]
\end{lstlisting}

Key: \texttt{comments} uses \texttt{operator.add} --- each node \textbf{appends} to the list rather than replacing it, enabling retry accumulation.

\subsection{Node Details}

\subsubsection{Analyze Node}
\begin{itemize}[leftmargin=1.5em]
    \item \textbf{Model:} \texttt{llama-3.3-70b-versatile}
    \item \textbf{Mode:} JSON response format
    \item \textbf{Temperature:} 0.3 (deterministic)
    \item \textbf{Detects:} post\_type (9 types), post\_tone (7 tones), sentiment, emotional\_tone, specific\_details, key\_insights, best\_response\_angles
    \item \textbf{Author context:} Injects post author name + headline into prompt
\end{itemize}

\subsubsection{Search Node (Conditional)}
\begin{itemize}[leftmargin=1.5em]
    \item \textbf{API:} SerpAPI Google Search
    \item \textbf{Triggered for:} advice\_tactical, hot\_take\_controversial, philosophical\_abstract, reflective\_lessons
    \item \textbf{Skipped for:} celebration\_win, milestone\_announcement, hiring\_team, self\_promo\_pitch
    \item \textbf{Query builder:} Uses topic + specific\_details + post\_type for focused queries
\end{itemize}

\subsubsection{RAG Node}
\begin{itemize}[leftmargin=1.5em]
    \item \textbf{Database:} 1,719 reference comments in SQLite
    \item \textbf{Query:} Filter by \texttt{post\_type}, sort by \texttt{engagement\_score DESC}, limit 5
    \item \textbf{Fallback:} Keyword search in \texttt{post\_content} if type match $<$ 3 results
    \item \textbf{Output:} Formatted examples + extracted angle names
\end{itemize}

\subsubsection{Generate Node (Parallel)}
\begin{itemize}[leftmargin=1.5em]
    \item \textbf{Model:} \texttt{llama-3.3-70b-versatile}
    \item \textbf{3 calls via \texttt{asyncio.gather()}} --- truly parallel
    \item \textbf{Angle selection priority:} RAG angles $\rightarrow$ AI-suggested valid angles $\rightarrow$ post type defaults
    \item \textbf{Only valid angles used:} 21 defined angles, rejects LLM-invented ones
    \item \textbf{Temperature escalation:} 0.70, 0.85, 1.00 for diversity
\end{itemize}

\subsubsection{Humanize + Validate Node}
\textbf{Humanization pipeline} (6 steps):
\begin{enumerate}[leftmargin=1.5em]
    \item Remove AI formality (``Furthermore'' $\rightarrow$ ``plus'', ``leverage'' $\rightarrow$ ``use'')
    \item Apply contractions (``do not'' $\rightarrow$ ``don't'', 35+ patterns)
    \item Add natural imperfections (drop trailing period, ellipsis, lowercase start)
    \item Trim to target length
\end{enumerate}

\textbf{Validation scoring:}
\begin{align*}
    \text{quality} &= 100 - \text{length\_penalty} - \text{ai\_phrase\_penalty} - \text{exclamation\_penalty}
\end{align*}

\textbf{Confidence formula:}
\begin{align*}
    \text{confidence} &= \underbrace{\frac{\text{quality}}{100} \times 0.40}_{\text{quality (40\%)}} + \underbrace{\text{angle\_bonus}}_{\text{30\%}} + \underbrace{0.15}_{\text{humanized}} + \underbrace{\text{word\_accuracy}}_{\text{10\%}} + \underbrace{\text{no\_issues}}_{\text{5\%}}
\end{align*}

Where angle\_bonus = 0.30 (RAG), 0.20 (good angle), or 0.10 (default).

\subsubsection{Retry Logic}
\begin{itemize}[leftmargin=1.5em]
    \item Triggered when any comment has quality $<$ 50
    \item Re-runs \texttt{generate} node with same state
    \item Maximum 2 retries
    \item After retries, top 3 comments by confidence are selected
\end{itemize}

% ============================================
% PROJECT STRUCTURE
% ============================================
\section{Project Structure}

\begin{lstlisting}
linkedin-comment-generator/
├── src/                        # Next.js Frontend
│   ├── app/
│   │   ├── page.tsx            # CSV upload page
│   │   ├── dashboard/page.tsx  # LinkedIn feed view
│   │   ├── post/[postId]/      # Post detail + comments
│   │   └── settings/page.tsx   # Commenter profile
│   ├── components/
│   │   ├── upload/             # CSV uploader, URL preview
│   │   ├── post/               # Post card (LinkedIn style)
│   │   └── comments/           # Comment card, actions
│   └── lib/                    # Shared types, data constants
│
├── backend/                    # FastAPI + LangGraph
│   ├── main.py                 # FastAPI app + CORS
│   ├── app/
│   │   ├── graph/
│   │   │   ├── state.py        # PipelineState TypedDict
│   │   │   ├── nodes.py        # 5 graph nodes
│   │   │   └── pipeline.py     # Graph compilation
│   │   ├── services/
│   │   │   ├── linkedin_scraper.py
│   │   │   ├── web_search.py
│   │   │   ├── rag_service.py
│   │   │   ├── prompt_engine.py
│   │   │   └── humanizer.py
│   │   ├── models/schemas.py   # Pydantic models
│   │   ├── routes/
│   │   │   ├── scrape.py       # /api/upload-csv, /api/scrape-profiles
│   │   │   ├── generate.py     # /api/generate-comments
│   │   │   └── posts.py        # /api/posts, /api/profile, /api/feedback
│   │   └── db/database.py      # SQLite connection + init
│   └── requirements.txt
│
├── data/                       # SQLite database
├── Dockerfile                  # Combined frontend + backend
├── start.sh                    # Starts both servers
├── next.config.ts              # Proxy /api/* -> FastAPI
└── docker-compose.yml
\end{lstlisting}

% ============================================
% DATABASE
% ============================================
\section{Database Schema}

\subsection{Tables}

\begin{center}
\small
\begin{tabularx}{\textwidth}{llX}
\toprule
\textbf{Table} & \textbf{Rows (typical)} & \textbf{Purpose} \\
\midrule
\texttt{scraped\_posts} & 10--1000 & Scraped LinkedIn posts with analysis \\
\texttt{generated\_comments} & 30--3000 & AI-generated comments (3 per post) \\
\texttt{comment\_feedback} & varies & User ratings for learning loop \\
\texttt{reference\_comments} & 1,719 & RAG dataset (high-engagement examples) \\
\texttt{commenter\_profiles} & 1--5 & Commenter voice profiles \\
\bottomrule
\end{tabularx}
\end{center}

\subsection{Key Fields: scraped\_posts}
\begin{lstlisting}
id, profile_url, author_name, author_headline,
author_profile_pic, post_text, post_url, post_urn,
likes_count, comments_count, shares_count,
media_type, media_urls (JSON), posted_at,
post_analysis (JSON), web_search_context (JSON),
status (scraped|analyzed|generated), batch_id
\end{lstlisting}

% ============================================
% API REFERENCE
% ============================================
\section{API Reference}

All endpoints are served by FastAPI on port 8000. Next.js proxies \texttt{/api/*} to FastAPI.

\begin{center}
\small
\begin{tabularx}{\textwidth}{llX}
\toprule
\textbf{Method} & \textbf{Endpoint} & \textbf{Description} \\
\midrule
POST & \texttt{/api/upload-csv} & Parse CSV, extract LinkedIn URLs \\
POST & \texttt{/api/scrape-profiles} & Scrape posts (with cache + dedup) \\
POST & \texttt{/api/generate-comments} & Full LangGraph pipeline \\
GET & \texttt{/api/posts} & List all scraped posts \\
GET & \texttt{/api/post/\{id\}} & Post detail + generated comments \\
GET & \texttt{/api/profile} & Get active commenter profile \\
POST & \texttt{/api/profile} & Save commenter profile \\
POST & \texttt{/api/feedback} & Rate a generated comment \\
POST & \texttt{/api/import-rag} & Import RAG CSV data \\
GET & \texttt{/api/rag-status} & Check RAG data count \\
GET & \texttt{/api/health} & Health check \\
\bottomrule
\end{tabularx}
\end{center}

% ============================================
% PROMPT ENGINEERING
% ============================================
\section{Prompt Engineering}

\subsection{Generation Prompt Structure}

The prompt sent to Groq for comment generation follows this structure:

\begin{lstlisting}
1. POST TO COMMENT ON (full text, 1500 chars max)
   + Author context (name, headline)
   + Post type, topic, key details

2. WHO YOU ARE (commenter identity)
   + Name, headline, expertise
   + Real comment examples (voice reference)
   + Signature phrases

3. COMMENT STRATEGY
   + Angle instruction (what to say)
   + Variation instruction (how to structure)

4. TONE MATCHING
   + Post tone detected -> matching instruction

5. RULES
   + Word count target
   + BANNED words/phrases (40+ AI tells)

6. RAG EXAMPLES (if available)
   + 5 high-engagement reference comments

7. WEB CONTEXT (if searched)
   + 3 search result snippets
\end{lstlisting}

\subsection{Anti-AI Detection}

The system actively avoids AI-detectable patterns:

\begin{itemize}[leftmargin=1.5em]
    \item \textbf{40+ banned words:} leverage, navigate, unlock, delve, foster, robust, seamless, holistic, innovative, transformative...
    \item \textbf{Banned phrases:} ``it's worth noting'', ``in today's'', ``couldn't agree more'', ``this resonates deeply''...
    \item \textbf{Structural:} No intro$\rightarrow$balanced$\rightarrow$summary$\rightarrow$conclusion pattern
    \item \textbf{Humanizer:} Contractions, imperfections, casual connectors
    \item \textbf{Tone matching:} Casual post $\rightarrow$ casual comment (no over-formality)
\end{itemize}

% ============================================
% COSTING
% ============================================
\section{Detailed Costing Analysis}

\subsection{API Pricing (as of April 2026)}

\begin{center}
\begin{tabularx}{\textwidth}{Xrr}
\toprule
\textbf{Service} & \textbf{Pricing} & \textbf{Free Tier} \\
\midrule
Groq llama-3.3-70b (input) & \$0.59 / M tokens & --- \\
Groq llama-3.3-70b (output) & \$0.79 / M tokens & --- \\
Groq llama-4-scout (input) & \$0.11 / M tokens & --- \\
Groq llama-4-scout (output) & \$0.34 / M tokens & --- \\
RapidAPI LinkedIn Scraper & \$0.025 / call & 10 calls/day \\
SerpAPI Google Search & \$0.015 / search & 100/month \\
Railway Hosting (Hobby) & \$5 / month & --- \\
\bottomrule
\end{tabularx}
\end{center}

\subsection{Per-Post Cost Breakdown}

For each post processed, the pipeline makes:

\begin{center}
\begin{tabularx}{\textwidth}{Xlrrr}
\toprule
\textbf{Step} & \textbf{Service} & \textbf{Calls} & \textbf{Tokens} & \textbf{Cost} \\
\midrule
Scrape post & RapidAPI & 1 & --- & \$0.0250 \\
Analyze post & Groq 3.3-70b & 1 & $\sim$1,400 & \$0.0014 \\
Image analysis & Groq 4-scout & 0.5 & $\sim$2,300 & \$0.0002 \\
Web search & SerpAPI & 0.5 & --- & \$0.0075 \\
Generate ×3 & Groq 3.3-70b & 3 & $\sim$5,400 & \$0.0042 \\
\midrule
\textbf{Total} & & \textbf{6} & & \textbf{\$0.038} \\
\bottomrule
\end{tabularx}
\end{center}

\subsection{Monthly Cost Projections}

\begin{center}
\renewcommand{\arraystretch}{1.3}
\begin{tabular}{|r|r|r|r|r|r|}
\hline
\textbf{Posts/Mo} & \textbf{Groq} & \textbf{RapidAPI} & \textbf{SerpAPI} & \textbf{Railway} & \textbf{Total} \\
\hline
100 & \$0.58 & Free tier & Free tier & \$5 & \textbf{\$5.58} \\
\hline
500 & \$2.90 & \$25 (Pro) & \$25 & \$5 & \textbf{\$57.90} \\
\hline
1,000 & \$5.80 & \$25 (Pro) & \$75 & \$5 & \textbf{\$110.80} \\
\hline
5,000 & \$29.00 & \$75 (Ultra) & \$150 & \$20 & \textbf{\$274.00} \\
\hline
10,000 & \$58.00 & \$150 (Mega) & \$275 & \$20 & \textbf{\$503.00} \\
\hline
\end{tabular}
\end{center}

\begin{archbox}{Cost Insight}
\textbf{Groq is negligible} --- even 10,000 posts/month costs only \$58 in AI inference. The dominant costs are \textbf{RapidAPI} (LinkedIn scraping) and \textbf{SerpAPI} (web search). Optimize by caching aggressively and skipping search for non-tactical posts.
\end{archbox}

\subsection{Cost Per Comment}

\begin{align*}
    \text{Cost per comment} &= \frac{\$0.038}{3} = \$0.013 \\
    \text{At 500 posts/month} &: \frac{\$57.90}{1500 \text{ comments}} = \$0.039 \text{ per comment (incl. subscriptions)}
\end{align*}

% ============================================
% DEPLOYMENT
% ============================================
\section{Deployment}

\subsection{Railway (Production)}

Single container runs both FastAPI and Next.js:

\begin{lstlisting}
# Dockerfile builds Next.js + installs Python deps
# start.sh launches both:
uvicorn main:app --port 8000 &    # Backend
node server.js                     # Frontend on $PORT
\end{lstlisting}

\textbf{Deploy:}
\begin{lstlisting}
railway up --detach --service linkedin-comment-gen
\end{lstlisting}

\textbf{Post-deploy:} Import RAG data (ephemeral storage resets on deploy):
\begin{lstlisting}
curl -X POST https://YOUR_URL/api/import-rag \
  -F "file=@rag_data.csv"
\end{lstlisting}

\subsection{Environment Variables}

\begin{center}
\begin{tabularx}{\textwidth}{lXl}
\toprule
\textbf{Variable} & \textbf{Description} & \textbf{Required} \\
\midrule
\texttt{GROQ\_API\_KEY} & Groq API key & Yes \\
\texttt{RAPIDAPI\_KEY} & RapidAPI key & Yes \\
\texttt{RAPIDAPI\_HOST} & RapidAPI host & Yes \\
\texttt{SERPAPI\_KEY} & SerpAPI key & Optional \\
\texttt{BACKEND\_URL} & FastAPI URL (default localhost:8000) & No \\
\texttt{PORT} & Frontend port (Railway sets this) & No \\
\bottomrule
\end{tabularx}
\end{center}

% ============================================
% FUTURE ENHANCEMENTS
% ============================================
\section{Future Enhancements}

\begin{enumerate}[leftmargin=1.5em]
    \item \textbf{Persistent storage} --- Railway Volume or external PostgreSQL for RAG data persistence
    \item \textbf{LangSmith integration} --- Visual traces of every pipeline run
    \item \textbf{Auto-post comments} --- Selenium/browser extension to post directly
    \item \textbf{Existing comments analysis} --- Scrape top comments to avoid repetition
    \item \textbf{Multi-user support} --- Multiple commenter profiles, team access
    \item \textbf{Scheduling} --- Auto-generate comments at optimal posting times
    \item \textbf{Fine-tuning} --- Once 50K+ comments collected, fine-tune model for brand voice
    \item \textbf{Video analysis} --- Full video transcription (currently only thumbnails)
    \item \textbf{Chrome extension} --- Generate comments directly on LinkedIn
\end{enumerate}

\end{document}
